\chapter{Budget}\label{ap:budget}

This appendix presents a detailed financial estimation for the project's development, covering material resources, human labour, and operational expenditures. The final budget is calculated using a structured framework divided into three main components:

\begin{itemize}
    \item \textbf{Equipment and Software Costs}: Calculated through depreciation, reflecting the exact timeframe these assets were utilized for the project.
    \item \textbf{Labour and Development Costs}: Evaluated based on the gross cost to the employer, rather than the employee's net salary.
    \item \textbf{Indirect Costs}: Set at a fixed rate of 15\% of the total combined equipment and personnel expenditures.
\end{itemize}

\section{Equipment and Software Costs} \label{subsec:budget:equipment}
The execution of this project required specific hardware resources, such as development workstations and testing devices, as well as software licenses for development environments and specialized tools. 

Rather than attributing the full purchase price of these assets to the project, a depreciation model is applied. This model determines the proportional hardware and software costs based strictly on the project's duration and the standard useful life of each asset. The depreciated value for each item is computed using the following linear depreciation formula:

\begin{equation}
    C_{project} = \frac{C_{total}}{useful\_life} \times months\_used
\end{equation}
\addequation{Equipment and software depreciation formula}

\subsection{Equipment Cost}
Table \ref{tab:equipment_cost} details the hardware components utilized during the project's development. A standard useful life of 60 months (5 years) has been applied to all physical equipment, calculating the proportional cost assigned to the two months of active project usage.

\begin{table}[H]
    \centering
    \caption{Equipment Cost}
    \label{tab:equipment_cost}
    \begin{tabular}{
        p{0.32\textwidth}
        >{\centering\arraybackslash}p{0.09\textwidth}
        >{\centering\arraybackslash}p{0.1\textwidth}
        >{\centering\arraybackslash}p{0.1\textwidth}
        >{\centering\arraybackslash}p{0.09\textwidth}
        >{\centering\arraybackslash}p{0.09\textwidth}
        >{\centering\arraybackslash}p{0.11\textwidth}}
        \toprule
        \textbf{Item} & \textbf{Amount} & \textbf{Unitary Cost (€)} & \textbf{Total Cost (€)} & \textbf{Useful Life (months)} & \textbf{Months Used} & \textbf{Assigned Cost (€)} \\
        \midrule
        Personal laptop & 1 & 1,000.00 & 1,000.00 & 60 & 2 & 33.33 \\
        NVIDIA Jetson Orin Nano \cite{NVIDIAJetsonAGXa} & 1 & 286.99 & 286.99 & 60 & 2 & 9.57 \\
        Raspberry Pi 5 Model B - 8GB \acrshort{RAM} \cite{ltdBuyRaspberryPib} & 3 & 89.95 & 269.85 & 60 & 2 & 9.00 \\
        \acrshort{PCIe} Peripheral Board Case & 3 & 9.33 & 27.99 & 60 & 2 & 0.93 \\
        27\acrshort{W} \acrshort{USB}-C Power Supply & 3 & 14.95 & 44.85 & 60 & 2 & 1.50 \\
        \acrshort{RPi} Active Cooler & 3 & 5.95 & 17.85 & 60 & 2 & 0.60 \\
        Micro \acrshort{SD} 64 GB & 3 & 9.95 & 29.85 & 60 & 2 & 1.00 \\
        Micro \acrshort{HDMI} cable & 1 & 7.82 & 7.82 & 60 & 2 & 0.26 \\
        \acrshort{AI} HAT+ \cite{AIHATsRaspberry} with Hailo-8 \cite{AIAcceleratorHailo8} & 1 & 124.95 & 124.95 & 60 & 2 & 4.17 \\
        \acrshort{AI} HAT+ \cite{AIHATsRaspberry} with Hailo-8L \cite{Hailo8LEntryLevelAIa} & 1 & 78.95 & 78.95 & 60 & 2 & 2.63 \\
        \acrshort{RPi} Camera Module 3 \cite{ltdBuyRaspberryPic} & 1 & 29.95 & 29.95 & 60 & 2 & 1.00 \\
        \acrshort{RPi} \acrshort{AI} Camera \cite{AICameraRaspberry} & 1 & 76.96 & 76.96 & 60 & 2 & 2.57 \\
        \acrshort{RPi}5 Camera cable & 1 & 3.00 & 3.00 & 60 & 2 & 0.10 \\
        \midrule
        \textbf{Total Equipment Cost} & & & \textbf{1,999.01} & & & \textbf{66.63} \\
        \bottomrule
    \end{tabular}
\end{table}

\subsection{Software Cost}
The software tools required for the design, development, and documentation phases are listed in Table \ref{tab:software_cost}. Fortunately, the majority of the environment relies on open-source or freely available tools, significantly mitigating licensing expenses. For commercial software, a standard 12-month useful life (representing an annual license) is applied.

\begin{table}[H]
    \centering
    \caption{Software Cost}
    \label{tab:software_cost}
    \begin{tabular}{
        p{0.25\textwidth}
        >{\centering\arraybackslash}p{0.1\textwidth}
        >{\centering\arraybackslash}p{0.1\textwidth}
        >{\centering\arraybackslash}p{0.1\textwidth}
        >{\centering\arraybackslash}p{0.1\textwidth}
        >{\centering\arraybackslash}p{0.1\textwidth}
        >{\centering\arraybackslash}p{0.11\textwidth}}
        \toprule
        \textbf{License} & \textbf{Amount} & \textbf{Unitary Cost (€)} & \textbf{Total Cost (€)} & \textbf{Useful Life (months)} & \textbf{Months Used} & \textbf{Assigned Cost (€)} \\
        \midrule
        Visual Studio Code & 1 & Free & Free & N/A & 2 & Free \\
        Docker & 1 & Free & Free & N/A & 2 & Free \\
        GitHub & 1 & Free & Free & N/A & 2 & Free \\
        Draw.io & 1 & Free & Free & N/A & 2 & Free \\
        Visual Paradigm \tablefootnote{ VP Standard \cite{EssentialUMLBPMN}.} & 1 & 1,100.50 & 1,100.50 & 12 & 2 & 183.42 \\
        Overleaf & 1 & Free & Free & N/A & 2 & Free \\
        Zotero & 1 & Free & Free & N/A & 2 & Free \\
        \midrule
        \textbf{Total Software Cost} & & & \textbf{1,100.50} & & & \textbf{183.42} \\
        \bottomrule
    \end{tabular}
\end{table}

\subsection{Total Equipment and Software Cost}
Table \ref{tab:total_equipment_software_cost} provides a consolidated view of the combined material resources.

\begin{table}[H]
    \centering
    \caption{Total Equipment and Software Cost}
    \label{tab:total_equipment_software_cost}
    \begin{tabular}{p{0.35\textwidth} p{0.25\textwidth}}
        \toprule
        \textbf{Category} & \textbf{Total Cost (€)} \\
        \midrule
        Total Equipment Cost & 66.63 \\
        Total Software Cost & 183.42 \\
        \midrule
        \textbf{Grand Total} & \textbf{250.05} \\
        \bottomrule
    \end{tabular}
\end{table}

By aggregating the depreciated costs of all individual hardware and software components utilized throughout the development lifecycle, the total equipment and software cost applied to this project is estimated at 250.05 €.

\section{Development and Labour Costs} \label{subsec:budget:labour}

Human resource expenses constitute the most significant portion of the project's budget. This section accounts for the total cost of employment, which extends beyond the base net salary to include social security contributions, taxes, and other mandatory employee benefits. The labour estimation is based on the following criteria:

\begin{itemize}
    \item The gross hourly labour cost for an engineer is established at 36.71 € (Value obtained from the \gls{INE} \cite{CosteLaboralPora}).
    \item The project required a total effort of 350 hours.
\end{itemize}

\vspace{-10pt}

\begin{equation}
    C_{labour} = hours \times C_{hour} = 350 \times 36.71 = 12,848.50\text{ €}
\end{equation}
\addequation{Total labour cost calculation}

\section{Indirect Costs} \label{sec:budget:indirect}
In addition to the direct expenses calculated in the previous sections, it is necessary to account for indirect operational costs. These encompass overhead expenses that are difficult to itemize per project but are essential for day-to-day operations. Examples include electricity consumption, high-speed Internet access, workspace maintenance, and general administrative overhead. 

Following standard financial estimation practices for engineering projects, these expenses are approximated as a 15\% surcharge applied to the sum of all direct material (equipment and software) and labour expenses:

\begin{equation}
    C_{indirect} = 0.15 \times (C_{equipment} + C_{labour}) = 0.15 \times (250.05 + 12,848.50) = 1,964.78\text{ €}
\end{equation}
\addequation{Total indirect cost calculation}

\section{Budget Summary} \label{sec:budget:summary}
A comprehensive breakdown of the estimated total project budget, aggregating direct and indirect costs, is illustrated in \autoref{tab:budget_summary}.

\begin{table}[H]
    \centering
    \caption{Budget Summary}
    \label{tab:budget_summary}
    \begin{tabular}{p{0.4\textwidth}p{0.2\textwidth}}
        \toprule
        \textbf{Item} & \textbf{Total Cost (€)} \\
        \midrule
        Equipment and Software & 250.05 \\
        Development and Labour & 12,848.50 \\
        Indirect Costs (15\%) & 1,964.78 \\
        \midrule
        \textbf{Total} & \textbf{15,063.34} \\
        \bottomrule
    \end{tabular}
\end{table}

\section{Conclusion} \label{sec:budget:conclusion}
The total projected budget for the development of the proposed system amounts to 15,063.34 €. As anticipated in technology and engineering projects, human resources constitute the vast majority of the expenses. However, by leveraging open-source software and applying a strict depreciation model for the hardware, material, and operational costs have been effectively minimised. This comprehensive estimate demonstrates the economic feasibility of the project, underlining its viability as a cost-effective and scalable solution for sustainable urban development and smart city initiatives.
